% ===================================================================
%  图表第 10 号 — 海。— 港口。
%  Traduction 2026 d'apres texto/io, controlee sur texto/fr, texto/en
%  et texto/es. Consigne : voir texto/zh/10-tubiao-01.tex.
%
%  Le francais decale sa numerotation d'un cran a partir de l'alinea 8.
%  C'est une coquille de l'imprimeur francais, conservee dans la
%  transcription ; le chinois suit l'ido, qui compte juste.
% ===================================================================

%%K t10-apar-1 apar
\VUsaut{4.0mm}
\VUtitre{15.0pt}{60}{图表第 10 号}
\VUsaut{3.0mm}
\VUpk{11.4pt}{40}{海。--- 港口。}
\VUsaut{3.0mm}

%%K t10-01-1 p
1. --- 夏天，有人到山里找清静和凉爽，也有人宁愿去海边。去年我们就是这样，
因为我们很爱\VUgras{海洋} \textsuperscript{(1)}。

%%K t10-02-1 p
2. --- 就我来说，望着那片一直铺到\VUgras{天边} \textsuperscript{(2)} 的无边水面，
看着巨大的\VUgras{轮船} \textsuperscript{(3)} 载着去美洲的旅客启程作长途
\VUgras{航行}，是极大的快乐。

%%K t10-03-1 p
3. --- 所以父亲知道我的喜好，放假总挑一处很安静的海滩，就在我们一个大
\VUgras{军港}附近。

%%K t10-03-2 p
上次我们住在那里时，\VUgras{舰队}来了，在那道巨大的\VUgras{防波堤}
\textsuperscript{(4)} 后面下了锚——防波堤封住并护着\VUgras{军港内湾}
\textsuperscript{(5)}——我
得以从容欣赏各样\VUgras{军舰}：从会下潜、会没入水中的小\VUgras{潜水艇}
\textsuperscript{(10)}，到那庞大的\VUgras{铁甲舰} \textsuperscript{(9)}——它至今除了
\VUgras{鱼雷艇} \textsuperscript{(8)} 的偷袭外，没有什么可怕的。

%%K t10-tit-2 sub
\VUsaut{3.0mm}
\VUpk{10.2pt}{40}{小艇。}
\VUsaut{3.0mm}

%%K t10-04-1 p
4. --- 鱼雷艇早就能很\VUgras{机灵}地找出厚铁甲的弱处，在那里安上危险的
\VUgras{鱼雷}，一炸就把船送到海底；可是它到底还是不如潜水艇可怕，因为
驱逐舰很容易追上它。

%%K t10-05-1 p
5. --- 我还看见了快速的装甲\VUgras{巡洋舰} \textsuperscript{(6)}，几乎同真正的
铁甲舰一样刀枪不入；还有几艘\VUgras{运兵船} \textsuperscript{(12)}、三四艘
\VUgras{炮舰} \textsuperscript{(7)}，最后还有一艘\VUgras{护卫舰} \textsuperscript{(11)}。
\VUgras{那支舰队起锚的样子最叫人看得有趣。}

%%K t10-06-1 p
6. --- 那些钢铁的庞然大物同商船队里还在用的漂亮\VUgras{三桅船}或秀气的
\VUgras{帆船} \textsuperscript{(13)}，样子相差多么大啊；我在\VUgras{码头堤}
\textsuperscript{(14)} 上散步时，常看见帆船张满帆驶进港来。当然，\VUgras{风}不
顺时它们就吃亏得多。进船坞要把帆全收起来，还得要一艘\VUgras{拖船}
\textsuperscript{(16)} 去把它们拖到岸边，像拖普通的\VUgras{驳船} \textsuperscript{(17)}
一样。

%%K t10-tit-3 sub
\VUsaut{3.0mm}
\VUpk{10.2pt}{40}{商港。}
\VUsaut{3.0mm}

%%K t10-07-1 p
7. --- 我说的这个港口所以有趣，是因为它不但有一座靠巨大工程人工造成的
军港内湾，还有一座极好的\VUgras{商港}，就在一条大河的\VUgras{河口}。

%%K t10-08-1 p
8. --- 隔开两个船坞的那条地带筑了工事，由一排伸进海里的\VUgras{炮台}和
\VUgras{棱堡}护着。要在\VUgras{城墙} \textsuperscript{(15)} 上散步得有特别的
许可证。我托叔叔很容易就拿到了一张，他在\VUgras{船厂} \textsuperscript{(18)}
当工程师；我便去看那些在大棚子底下造着的船。

%%K t10-09-1 p
9. --- 我甚至看见了一艘铁甲舰下水，还记得那庞然大物一放开就顺着船台滑
下去、浮在河水上时，全场人心里那阵激动。

%%K t10-10-1 p
10. --- 要到船厂去，得穿过\VUgras{操场} \textsuperscript{(19)}——驻军每天在那里
操练——再走过一座铁\VUgras{便桥} \textsuperscript{(20)}，
它连着操场和\VUgras{军械库} \textsuperscript{(21)}。

%%K t10-tit-4 sub
\VUsaut{3.0mm}
\VUpk{10.2pt}{40}{军械库和船坞。}
\VUsaut{3.0mm}

%%K t10-11-1 p
11. --- 我跟叔叔参观了那座大厂，里面满是\VUgras{炮} \textsuperscript{(22)}、
\VUgras{炮弹} \textsuperscript{(23)} 和\VUgras{开花弹}。我还看了\VUgras{锻工场} \textsuperscript{(24)}，
轮船的\VUgras{锅炉} \textsuperscript{(25)} 就是在那里造的。

%%K t10-12-1 p
12. --- 近旁就是\VUgras{船坞} \textsuperscript{(26)}——干船坞——还有很值得看的
\VUgras{浮船坞} \textsuperscript{(27)}；在一艘正在修的船上，我近看了船的
\VUgras{螺旋桨} \textsuperscript{(28)}、\VUgras{龙骨} \textsuperscript{(29)} 和
\VUgras{舵} \textsuperscript{(30)}。

%%K t10-13-1 p
13. --- 商港同军港一样值得琢磨，我在码头上消磨了好些钟头，那里泊着上溯
大河的\VUgras{明轮船} \textsuperscript{(31)}；我看着\VUgras{人字起重机}
\textsuperscript{(32)} 干活，把极重的东西提起来，像提羽毛似的。

%%K t10-tit-5 sub
\VUsaut{4.0mm}
\VUpk{10.2pt}{40}{领港。}
\VUsaut{3.0mm}

%%K t10-14-1 p
14. --- 我很欣赏那些\VUgras{远洋邮船} \textsuperscript{(34)} 打着\VUgras{旗}
\textsuperscript{(35)} 驶出锚地，由一位高明的\VUgras{领港} \textsuperscript{(36)} 掌
着，他极会顺着用\VUgras{浮标} \textsuperscript{(40)} 和\VUgras{标杆}标出的
\VUgras{航道}走，躲开那些只靠一张方帆、极难驾驭的\VUgras{驳运船}
\textsuperscript{(41)}。我看得出\VUgras{船头} \textsuperscript{(37)} 是二等
\VUgras{舱} \textsuperscript{(39)}，\VUgras{船尾} \textsuperscript{(38)} 是头等客人的
大厅。

%%K t10-15-1 p
15. --- 有时我也看\VUgras{货船} \textsuperscript{(42)} 装货，把从\VUgras{货栈}
\textsuperscript{(43)} 和\VUgras{港站} \textsuperscript{(44)} 运来的货堆进去，那些货
是\VUgras{码头铁道} \textsuperscript{(45)} 上一列列火车送来的。

%%K t10-tit-6 sub
\VUsaut{3.0mm}
\VUpk{10.2pt}{40}{划船消遣。}
\VUsaut{3.0mm}

%%K t10-16-1 p
16. --- 我常在\VUgras{内船坞} \textsuperscript{(53)} 里划船，小\VUgras{艇}
\textsuperscript{(46)} 都躲在那里；握着\VUgras{桨} \textsuperscript{(47)} 是我的一桩
乐事。我会爬上\VUgras{甲板} \textsuperscript{(48)}——那是一条
\VUgras{帆船} \textsuperscript{(49)} 的甲板——顺着\VUgras{索具} \textsuperscript{(52)}
攀上\VUgras{桅} \textsuperscript{(51)}，到\VUgras{帆桁} \textsuperscript{(50)} 上
爬，然后再坐一条
\VUgras{小艇} \textsuperscript{(54)} 在沉重的\VUgras{渔船} \textsuperscript{(55)} 中间
接着游；渔船上挂着\VUgras{渔网} \textsuperscript{(56)} 晾着。

%%K t10-17-1 p
17. --- \VUgras{码头} \textsuperscript{(58)} 上，最形形色色的\VUgras{货物}
\textsuperscript{(59)} 从我眼前过去，装在\VUgras{木箱} \textsuperscript{(60)} 里或
打成\VUgras{包} \textsuperscript{(61)}。它们就是驳船的\VUgras{载货}
\textsuperscript{(63)}；有力的\VUgras{起重机} \textsuperscript{(62)} 把它们吊上来，
放到\VUgras{货车} \textsuperscript{(64)} 上，同时\VUgras{承运人}在\VUgras{海关}
\textsuperscript{(65)} 报关纳税。

%%K t10-18-1 p
18. --- 最后，走到\VUgras{开合桥} \textsuperscript{(66)} 或\VUgras{船闸}
\textsuperscript{(69)}，经过那条\VUgras{挖泥船} \textsuperscript{(68)}，
它正在清\VUgras{运河} \textsuperscript{(67)}，我就在\VUgras{信号杆} \textsuperscript{(70)}
旁边靠上堤岸。

%%K t10-19-1 p
19. --- 就在那道堤的另一边，把舰队军官送上岸的\VUgras{汽艇} \textsuperscript{(71)}
靠拢过来。看水兵们合着拍子举起桨，小艇随着\VUgras{浪} \textsuperscript{(72)}
轻轻靠上台阶，是很好看的。在这里最容易看清\VUgras{潮汐}，看清
\VUgras{沙滩} \textsuperscript{(73)} 上\VUgras{退潮}和\VUgras{涨潮}的来去。

%%K t10-tit-7 sub
\VUsaut{3.0mm}
\VUpk{10.2pt}{40}{军官俱乐部。}
\VUsaut{3.0mm}

%%K t10-20-1 p
20. --- 回家我坐\VUgras{电车} \textsuperscript{(74)}，它的\VUgras{站}
\textsuperscript{(75)} 就在军官俱乐部门前那根\VUgras{柱子}那儿。

%%K t10-20-2 p
从俱乐部的\VUgras{阳台} \textsuperscript{(76)} 上——阳台装点着\VUgras{锚}
\textsuperscript{(77)}、炮和炮弹——我常看见一大群海军军官：佩剑的\VUgras{中尉}
\textsuperscript{(78)}，佩\VUgras{军刀}的炮兵\VUgras{上尉} \textsuperscript{(79)} 和
\VUgras{步兵} \textsuperscript{(80)} 上尉。他们身后站着一个\VUgras{水兵}
\textsuperscript{(81)}，他们要看远处的动静时，他就递上\VUgras{望远镜}。

%%K t10-21-1 p
21. --- 我也见过年轻的\VUgras{少尉} \textsuperscript{(82)} 从\VUgras{舰长}
\textsuperscript{(83)} 或\VUgras{将军} \textsuperscript{(84)} 那里领命令，
\VUgras{军需长} \textsuperscript{(85)} 在一旁等着把命令带回船上。

%%K t10-22-1 p
22. --- 从那阳台上望出去景色极好：整个海滩尽收眼底，滩上有\VUgras{帐篷}
\textsuperscript{(86)} 和\VUgras{更衣棚} \textsuperscript{(87)}；到洗海水澡的时候，
看得见\VUgras{洗海澡的人} \textsuperscript{(88)} 在\VUgras{赌场} \textsuperscript{(89)}
前面快活地嬉戏。

%%K t10-23-1 p
23. --- 海滩尽头，岸线伸成一个小小的石头\VUgras{半岛} \textsuperscript{(91)}，
\VUgras{碎浪} \textsuperscript{(92)} 冲在上面；半岛上筑着一座\VUgras{灯塔}
\textsuperscript{(93)}，夜里给水手指路。这半岛只靠一条很窄的\VUgras{地峡}
\textsuperscript{(90)} 连着陆地。

%%K t10-tit-8 sub
\VUsaut{3.0mm}
\VUpk{10.2pt}{40}{海岸的全景。}
\VUsaut{3.0mm}

%%K t10-24-1 p
24. --- \VUgras{断崖} \textsuperscript{(94)} 一直延伸下去，被海凿开，海在它上面
刻出一串奇形怪状的\VUgras{海湾} \textsuperscript{(95)} 和\VUgras{海角}，使它
分外好看。

%%K t10-24-2 p
\VUgras{高地} \textsuperscript{(96)} 俯瞰着\VUgras{海峡}
\textsuperscript{(98)}，上面有一座很要紧的\VUgras{炮台} \textsuperscript{(97)} 和
几幢别墅。

%%K t10-25-1 p
25. --- 那道海峡把一座很危险的\VUgras{岛} \textsuperscript{(99)} 同大陆隔开，
岛四周是\VUgras{礁石} \textsuperscript{(100)} 和\VUgras{暗礁}，小船甚至大船
有时会在那里失事。尽管救援快、\VUgras{救生艇}到得及时，这些\VUgras{海难}
仍然可怕；在分点前后的\VUgras{大风}里，好几条船连人带货都沉了。

%%K t10-25-2 p
虽有这样的近邻，这座\VUgras{港口}水手们还是常来。
\VUsaut{6.0mm}
\VUfilet{22mm}
